\documentclass{article}
% Shared layout. Keep benchmark prose and result-specific plots outside this folder.
\PassOptionsToPackage{numbers}{natbib}
\usepackage[preprint]{neurips_2026}
\usepackage[T1]{fontenc}
\usepackage[hidelinks]{hyperref}
\usepackage{booktabs,amsmath,microtype,pgfplots}
\pgfplotsset{compat=1.17}
\usetikzlibrary{patterns}
\definecolor{okblue}{HTML}{0072B2}
\definecolor{okorange}{HTML}{D55E00}
\pgfplotsset{every axis/.append style={font=\small, tick label style={font=\small}, width=\linewidth, height=5cm, grid=major}}

\newcommand{\PaperTitle}{Tool Errors: Measuring Safe Recovery from Ambiguous API Failures}
\newcommand{\PaperAbstract}{An API error does not establish whether an operation took effect. Tool Errors measures safe recovery through stateful fault injection and replay-based scoring. Its 40 public episodes cover 6 fault families across 4 domains. We execute 4960 episodes using scripted controls and a simulated status-read omission sweep. The reference and response-aware controller each reach full credit on 40/40 tasks (95\% Wilson interval 91.2-100.0\%). Removing reads in a simulated ablation leaves 7/40 full-credit tasks (8.7-31.9\%). At simulated read-omission probability 0.5, the mean full-credit fraction is 53.8\% (50.6-56.9\%, seed bootstrap). These are checks of the evaluation design, not trained-model results. No hosted model, human study, or production deployment is evaluated.}
\newcommand{\AuthorOrcid}{https://orcid.org/0009-0005-9964-6841}
\newcommand{\NTasks}{40}
\newcommand{\NFaults}{6}
\newcommand{\NDomains}{4}
\newcommand{\NEpisodes}{4960}
\newcommand{\NConditions}{12}
\newcommand{\NSeeds}{20}
\newcommand{\NLevels}{5}
\newcommand{\CiLevel}{95}
\newcommand{\NBoot}{2000}
\newcommand{\NCalls}{12}
\newcommand{\NTurns}{12}
\newcommand{\FullGood}{100.0 [91.2, 100.0]\%}
\newcommand{\FullNoRead}{17.5 [8.7, 31.9]\%}
\newcommand{\MidSweep}{53.8 [50.6, 56.9]\%}
\newcommand{\NoReadCount}{7}
\newcommand{\RateCount}{4}
\newcommand{\Recoverable}{33}

\title{\PaperTitle}
\author{Pooria Arab\\Independent researcher\\\href{\AuthorOrcid}{\AuthorOrcid}}
\begin{document}
\maketitle
\begin{abstract}\PaperAbstract\end{abstract}
\section{Why this matters}
A failed response does not necessarily mean a failed operation. An invoice may exist even when its creation call times out. An email may reach its recipient before a connection closes. Repeating the request can therefore produce an unwanted second effect. Stopping immediately can leave the requested work unfinished. Reporting success can mislead the user when the actual outcome remains unknown.\par
Tool Errors measures safe, bounded recovery from tool faults. Its target behavior combines useful action with evidence about the resulting state. The agent must preserve the requested operation identity, respond to the observed error, and report an outcome supported by the trace. The construct is narrower than general tool use. It concerns what happens after a tool fails or returns an incomplete result.\par
The benchmark contributes a controlled environment for separating these behaviors. It supplies fault schedules, observable status queries, and an outcome report action. Scoring replays the submitted actions in a fresh environment. This makes the outcome independent of a submitted success claim. The resulting evidence concerns a defined recovery contract, rather than an evaluator's impression of the final answer.\par
This paper studies the measurement design without presenting a learned-model comparison. Scripted positive and negative controls test score separation. Simulated control ablations test which recovery safeguards affect the score. A further simulated population experiment varies the probability of omitting a status read. These experiments establish reproducible behavior within the simulator. They do not establish the prevalence of these errors in deployed agents.\par
\subsection{A timeout creates an information problem}
Consider a write followed by a timeout. There are two relevant possibilities: the server rejected the operation, or the server completed it without returning a response. These possibilities demand different next actions. Retrying is useful in the first case. Reading the operation status can avoid a duplicate in the second case. The benchmark makes the second case explicit through a committed-timeout fault.\par
The important distinction is between attempting a call and establishing its effect. A model can produce a plausible recovery explanation without checking the actual state. Conversely, a controller can complete the operation but leave the user without a reliable report. The evaluation therefore retains separate evidence for completion, side-effect safety, retry discipline, and reporting. A single mean remains available, but the components explain why it changes.\par
\section{Related work}
Tool-agent evaluation already extends beyond static question answering. The $\tau$-bench environment combines domain tools, user interaction, and policy constraints. It evaluates final database state and repeated-run reliability \citep{yao2024tau}. Tool Errors isolates a narrower problem: recovery under explicit fault schedules. This focus complements broader interaction tasks. We do not claim that tool faults are absent from every existing benchmark.\par
The Agentic Benchmark Checklist identifies reward and task-design defects as threats to agent evaluation \citep{zhu2025abc}. It motivates our use of positive controls, trivial policies, and replayed outcomes. BetterBench similarly treats benchmark design and documentation as objects of assessment \citep{reuel2024betterbench}. Our contribution is not a new checklist. It is a concrete recovery environment accompanied by evidence about its scoring behavior.\par
Construct validity asks whether a measure represents the phenomenon named by its claims \citep{bean2025construct}. Here, the relevant bridge is explicit. A request identifies an intended effect; the trace records attempted actions; the environment tracks effects; and the score evaluates recovery against that state. This bridge supports claims about the defined simulator. Transfer to real services still requires evidence about realistic fault distributions, status semantics, and operational consequences.\par
PayBench frames delegated action through useful completion, constraints, and control conditions \citep{plunkett2026paybench}. We follow its construct-first presentation. Tool Errors does not inherit its matched trap/lookalike design. Our comparisons pair policies on identical fault-bearing tasks. There is no no-fault matched partner in this study, so the results cannot identify a causal penalty for introducing faults into otherwise identical tasks.\par
\section{Experiment design}
\subsection{Episode contract}
Each episode requests one operation in invoicing, email, calendar, or inventory. The visible instructions provide an operation key and requested value. A write tool exposes a payload field. Additional tools read operation status, advance a virtual clock, and report the outcome. The operation key identifies the requested effect across retries. Preserving that identity is part of the task, not an optional optimization.\par
The environment follows a fixed fault schedule. Tool responses can indicate a rate limit, a temporary server failure, a committed timeout, a changed payload field, partial completion, or permanent failure. These are synthetic conditions with controlled state transitions. They are not recordings of vendor incidents. The virtual clock allows waiting policies to be tested without physical delays.\par
The agent has access to instructions, tool descriptions, and preceding observations. The response-aware controller in this study uses only those visible fields. It does not inspect the task's fault label. The packaged reference is different: it selects a complete trace using privileged task metadata. Its role is to demonstrate reachability of full credit. It is an oracle-style control, not a model baseline.\par
\subsection{Controls and hypotheses}
The primary positive controls are the packaged reference and a response-aware scripted controller. The latter parses the visible request, sends the documented payload, and branches on tool results. It waits after a rate limit, retries a temporary server error, follows a revealed field rename, and reads after ambiguous or completed writes. It reports failure after a permanent error. That report follows an attempt rather than replacing the requested work with a refusal.\par
Negative controls produce an empty answer, a constant answer, an echoed request, a refusal, or repeated calls to the first tool with empty arguments. A separate false-success policy reports success without acting. These policies test distinct shortcuts. Plain prose does not submit an outcome report. Repeated invalid calls consume the bounded call budget. A success claim has no authority to establish an effect.\par
The simulated ablations remove all waits, all status reads, or the final report from the privileged reference. A further simulated condition inserts a redundant write before the report. These changes preserve the requested values and operation keys. They target control behavior rather than language understanding. Removing a control can affect subsequent state, so these are total trace interventions, not isolated changes to a score component.\par
Before data collection, we hypothesized that both positive controls would reach full credit and the trivial policies would not. We also expected each simulated ablation to reduce the mean. The additional simulated sweep tests whether increasing read omissions reduces full credit. These hypotheses describe diagnostic separation. We do not conduct a learned-model ranking or select a statistical significance threshold.\par
\subsection{Execution and study ordering}
The study uses \NConditions{} fixed conditions on \NTasks{} public episodes. Each condition runs twice. Repetition checks deterministic semantic traces and does not enlarge the sample. The additional simulated sweep uses \NLevels{} omission settings and \NSeeds{} seeds. Across both studies, the runner executes \NEpisodes{} episodes. All experiments use local scripted adapters and incur no model API charges.\par
The plan entered version control before study output was saved. This establishes local ordering, not independent preregistration. Two aborted adapter runs required repairs: successful results were parsed as wrapped objects, and a silent script lacked a terminal empty response. Both changes were logged before rerunning the unchanged grid. Infrastructure errors stop collection; they do not become benchmark failures or excluded tasks.\par
The runner permits \NTurns{} turns and \NCalls{} tool calls per episode, uses serial execution, and disables infrastructure retries. An episode can end earlier when the outcome report is recorded. The fault simulator still imposes its own stricter write limits. Thus, a runner budget does not authorize arbitrary retries. The controller and environment enforce different constraints.\par
\section{Coverage}
The public suite contains \NTasks{} synthetic episodes across \NDomains{} domains and \NFaults{} fault families. Table~\ref{tab:coverage} gives the actual cross-tabulation. Every domain appears under every fault family, but fault counts differ. The domain labels broaden the surface tasks while retaining a shared recovery contract. They do not establish equal realism across business workflows.\par
\begin{table}[ht]\centering\begin{tabular}{lrrrrr}\toprule
Fault & Invoice & Email & Calendar & Inventory & Total\\ \midrule
Schema drift & 4 & 1 & 1 & 3 & 9\\
Committed timeout & 2 & 2 & 1 & 2 & 7\\
Partial success & 1 & 2 & 2 & 1 & 6\\
Permanent failure & 1 & 2 & 2 & 2 & 7\\
Transient server error & 1 & 2 & 3 & 1 & 7\\
Rate limit & 1 & 1 & 1 & 1 & 4\\
\bottomrule\end{tabular}
\caption{Observed public task inventory. These are fixed counts, not estimated rates.}\label{tab:coverage}\end{table}
\begin{figure}[ht]\centering
\begin{tikzpicture}\begin{axis}[ybar,bar width=20pt,ymin=0,ymax=11,ylabel={Episode count},xtick={0,1,2,3,4,5},xticklabels={Schema,Timeout,Partial,Permanent,Server,Rate limit}]
\addplot[fill=okblue,nodes near coords] table[x=x,y=count,col sep=comma]{figures/coverage.generated.csv};
\end{axis}\end{tikzpicture}
\caption{Coverage includes all fault families but assigns different counts. The fixed suite contains \NTasks{} episodes. Counts describe the inventory, not estimated deployment frequencies.}\label{fig:coverage}\end{figure}

A full-credit aggregate weights each public task equally. It therefore inherits the observed fault mixture. Removing waits affects only rate-limit tasks in this design, while removing reads affects most recoverable operations. Comparing those aggregate penalties does not measure the relative importance of waiting and verification in production. It measures their exposure under this suite's composition.\par
Task generation varies operation identities, values, and payload fields. These variations discourage simple reuse of a single complete trace. They do not create new recovery semantics. Public task visibility also permits specialized controllers to recognize the environment. The present study intentionally uses such controllers to test the measurement path. A generalization study needs separately held-out instances and broader fault combinations.\par
Coverage remains limited in several ways. Faults occur in isolated episodes rather than a shared service with competing users. Status reads are reliable and current. Partial completion uses a controlled primary/secondary pattern. The environment does not model delayed replication, authorization changes, time-varying quotas, or irreversible external messages beyond its simulated state. These omissions bound the construct represented by the benchmark.\par
\section{Taxonomy of recovery failures}
\subsection{Timing errors and bounded retries}
A rate limit communicates when a retry becomes acceptable. A controller that immediately repeats the write violates the temporal contract even if its payload and key are correct. The virtual clock separates compliance with this instruction from elapsed wall time. A temporary server error tests a simpler response: retry the same intended operation after the initial failed attempt.\par
Permanent failure changes the goal. The desired state becomes an attempted operation with no effect and an accurate failure report. Persistent retrying does not improve that state. The benchmark limits writes and rejects an otherwise plausible report when attempts exceed that bound. This defines a useful stopping behavior without requiring the agent to infer whether a real service will eventually recover.\par
\subsection{Ambiguity, partial work, and duplicate effects}
A committed timeout leaves an effect despite the error response. Reading by operation key resolves the ambiguity. A second write before that read violates verification discipline. A later redundant write after known completion is also rejected. This is intentionally stricter than counting distinct stored records: redundant attempts can violate the task contract even when the simulator deduplicates storage.\par
Partial success requires a different response. A status read establishes that the requested operation is incomplete. The controller then resumes under the same key and verifies the resulting completion. Treating a partial result as success skips necessary work. Treating it as a total failure can also prompt an unsafe new operation. The shared key and intervening read distinguish recovery from a fresh request.\par
\subsection{Schema adaptation and unsupported reports}
A schema-drift task initially documents an obsolete payload field. The failed write reveals the replacement. Full-credit recovery uses that evidence to repair the next request. Sending the replacement field before the reveal is treated as guessing. This rule prevents privileged knowledge of the live field from masquerading as observed adaptation. It also limits applicability to systems where discovery through an error is the intended workflow.\par
Outcome reporting is grounded in replayed state. A success report must follow a status read of the written key after the last write. A truthful failure report must describe the failed operation. The message check is lexical, while the status check is state based. Neither a fluent final answer nor an externally supplied state snapshot replaces these checks.\par
The taxonomy separates absence of action from safe recovery. A refusal avoids duplicate effects but does not attempt the authorized task. A silent successful trace completes work without reporting it. A false-success trace supplies a report without supporting work. These distinctions explain why no single component can stand in for the complete recovery contract.\par
\section{Metrics and scoring}
\subsection{Replay and component credit}
Scoring reconstructs state by replaying tool names and arguments in a fresh seeded environment. Recorded result claims do not supply ground truth. Let $E$, $S$, $R$, and $H$ denote binary indicators for the expected end state, side-effect safety, retry discipline, and an honest report. The official score combines these indicators, then applies safety gates and caps. The weights and thresholds are stated in the accompanying generated scoring specification.\par
The base score is $(45E+25S+20R+10H)/100$. Recoverable tasks permit at most 4 writes; permanent failures permit 3. Unsupported success or failure claims cap the score at 0.05. Missing reports, unverified retries, and guessed schemas cap it at 0.35. Duplicate effects and excess writes force zero.

For recoverable faults, end-state credit requires the intended completed effect. For permanent failure, it requires at least one attempt and no effect. Safety and retry credit depend on reaching that expected state. Inaction therefore cannot accumulate those components on an ordinary recovery task. Reporting credit separately requires a matching status and a nonempty, suitable message.\par
The gates matter more than a small shift in component weights. A duplicate or excessive retry count forces the score to zero. Unsupported outcome claims receive a low cap. Missing reports and unverified retries receive another cap. As a result, the reported mean is a utility-like diagnostic score, not the probability of a correct episode. We report full-credit frequency separately to prevent that interpretation.\par
\subsection{Uncertainty and paired comparisons}
Binary frequencies carry \CiLevel\% Wilson intervals. These intervals describe uncertainty under a hypothetical sampling interpretation of the task set. The actual public tasks form a fixed synthetic inventory, not an independently sampled deployment population. Small fault strata have wide intervals, even when every observed task receives the same outcome. An observed ceiling is therefore not evidence of universal success.\par
Mean scores and paired score changes use \NBoot{} bootstrap draws stratified by fault family. Each draw preserves the fault counts while sampling tasks within each family. Comparisons use the same task in both conditions before resampling. This pairing removes task-composition differences from the contrast. The intervention remains specific to the reference trace and cannot identify its effect on an arbitrary learned policy.\par
Several intervals collapse because scores are constant within each fault family. This is a property of the conditional bootstrap, not evidence that real-world uncertainty vanishes. Wilson intervals on full-credit frequency remain visible alongside the collapsed mean intervals. The paired response-aware versus reference score contrast is null. It establishes equality on this inventory rather than equivalence across unseen environments.\par
\section{Results}
\subsection{Positive and negative controls}
Table~\ref{tab:results} reports every planned condition. Both positive controls reach full credit, with task-level frequency \FullGood{}. The privileged reference establishes that full recovery is reachable. The response-aware controller establishes that the visible observations suffice for its scripted decision rule. Neither result estimates the capability of a trained language model.\par
\begin{table}[ht]\centering\small\begin{tabular}{lrr}\toprule
Condition & Mean score [CI] & Full credit, \% [CI]\\ \midrule
Reference (privileged) & 1.000 [1.000, 1.000] & 100.0 [91.2, 100.0]\\
False success & 0.000 [0.000, 0.000] & 0.0 [0.0, 8.8]\\
Response-aware recovery & 1.000 [1.000, 1.000] & 100.0 [91.2, 100.0]\\
Empty & 0.000 [0.000, 0.000] & 0.0 [0.0, 8.8]\\
Constant & 0.000 [0.000, 0.000] & 0.0 [0.0, 8.8]\\
Echo & 0.000 [0.000, 0.000] & 0.0 [0.0, 8.8]\\
Refusal & 0.000 [0.000, 0.000] & 0.0 [0.0, 8.8]\\
First tool & 0.000 [0.000, 0.000] & 0.0 [0.0, 8.8]\\
Simulated: no wait & 0.900 [0.900, 0.900] & 90.0 [76.9, 96.0]\\
Simulated: no read & 0.216 [0.216, 0.216] & 17.5 [8.7, 31.9]\\
Simulated: no report & 0.350 [0.350, 0.350] & 0.0 [0.0, 8.8]\\
Simulated: redundant write & 0.175 [0.175, 0.175] & 17.5 [8.7, 31.9]\\
\bottomrule\end{tabular}
\caption{All planned controls. Brackets give \CiLevel\% intervals: fault-stratified bootstrap for mean scores and Wilson for full credit. Simulated conditions mutate reference traces.}\label{tab:results}\end{table}
The empty, constant, echo, refusal, first-tool, and false-success controls receive no full-credit tasks. Their exact intervals appear in Table~\ref{tab:results}. This evidence rejects those particular shortcuts. It does not prove that the grader resists every fixed sequence or adversarial submission. The controls differ deliberately in action and reporting behavior, so their shared floor should not be read as behavioral equivalence.\par
The refusal control produces prose without a report action. Its result therefore tests whether that prose receives recovery credit, not whether every possible failure report scores zero. Similarly, the first-tool control sends empty arguments. Its poor result does not show that all non-adaptive scripts fail. These narrow interpretations preserve the link between the implemented control and the claim it supports.\par
\subsection{Simulated control ablations}
Removing waits damages rate-limit tasks while leaving other fault families unchanged. Removing reads leaves only \NoReadCount{} full-credit tasks, giving \FullNoRead{}. Permanent failures remain fully credited because the contract does not require reading a nonexistent successful effect. The simulated no-read condition is therefore not a uniform degradation across faults.\par
\begin{figure}[ht]\centering
\begin{tikzpicture}\begin{axis}[ybar,bar width=20pt,ymin=0,ymax=1.12,ylabel={Full-credit fraction},xtick={0,1,2,3,4,5},xticklabels={Reference,Reactive,No wait,No read,No report,Extra write},x tick label style={rotate=20,anchor=east}]
\addplot[fill=okblue,error bars/.cd,y dir=both,y explicit] table[x=x,y=value,y error minus=low,y error plus=high,col sep=comma]{figures/controls.generated.csv};
\end{axis}\end{tikzpicture}
\caption{Simulated control ablations separate completion from verified recovery. Each condition uses the same \NTasks{} tasks. Error bars show \CiLevel\% Wilson intervals over tasks. The last four conditions are simulated mutations of reference traces.}\label{fig:controls}\end{figure}

Removing reports preserves the completed operations but triggers the missing-report cap. Its full-credit frequency falls to the floor. Adding a redundant write rejects recoverable operations that already reached known completion. Permanent failures remain within their attempt bound in this intervention. These outcomes show that the score responds to trace discipline as well as final completion.\par
\begin{table}[ht]\centering\small\begin{tabular}{lrr}\toprule
Simulated control removed/added & Paired score change [CI] & Full-credit count\\ \midrule
Simulated: no wait & -0.100 [-0.100, -0.100] & 36/40\\
Simulated: no read & -0.784 [-0.784, -0.784] & 7/40\\
Simulated: no report & -0.650 [-0.650, -0.650] & 0/40\\
Simulated: redundant write & -0.825 [-0.825, -0.825] & 7/40\\
\bottomrule\end{tabular}
\caption{Simulated paired ablation effects relative to the privileged reference. Bootstrap intervals preserve fault composition. Collapsed intervals reflect constant within-family scores.}\label{tab:contrasts}\end{table}
The simulated ablations support the directional hypotheses on this task inventory. They also expose planned nulls. Waiting removal has no effect outside rate limits. Read removal and the added write leave permanent-failure full credit unchanged. Reporting the nulls prevents the aggregate penalty from implying a universal role for every safeguard.\par
\begin{table}[ht]\centering\scriptsize\begin{tabular}{lrrr}\toprule
Fault & No wait, \% [CI] & No read, \% [CI] & Redundant write, \% [CI]\\ \midrule
Schema drift & 100.0 [70.1, 100.0] & 0.0 [0.0, 29.9] & 0.0 [0.0, 29.9]\\
Committed timeout & 100.0 [64.6, 100.0] & 0.0 [0.0, 35.4] & 0.0 [0.0, 35.4]\\
Partial success & 100.0 [61.0, 100.0] & 0.0 [0.0, 39.0] & 0.0 [0.0, 39.0]\\
Permanent failure & 100.0 [64.6, 100.0] & 100.0 [64.6, 100.0] & 100.0 [64.6, 100.0]\\
Transient server error & 100.0 [64.6, 100.0] & 0.0 [0.0, 35.4] & 0.0 [0.0, 35.4]\\
Rate limit & 0.0 [0.0, 49.0] & 0.0 [0.0, 49.0] & 0.0 [0.0, 49.0]\\
\bottomrule\end{tabular}
\caption{Fault-specific full-credit frequencies for simulated ablations, with \CiLevel\% Wilson intervals. Each column uses identical tasks.}\label{tab:faults}\end{table}
\subsection{Additional virtual experiment: simulated population sweep}
The simulated population experiment removes each reference status read with a specified probability. A hash of simulation seed, task identifier, and read position produces the draw. The same draws apply across omission settings. This common-random-number design supports a clean ordered intervention without changing operation values, error schedules, or report wording.\par
Each simulated seed yields a complete suite mean and a full-credit fraction. We average those seed-level summaries and bootstrap seeds. We do not pool every episode as an independent new task. The resulting intervals quantify variation induced by the simulated read-omission process. They exclude uncertainty about task selection, model behavior, and whether the omission mechanism resembles real agents.\par
\begin{figure}[ht]\centering
\begin{tikzpicture}\begin{axis}[xmin=0,xmax=1,ymin=0,ymax=1.08,xlabel={Simulated probability of omitting each status read},ylabel={Full-credit fraction}]
\addplot[color=okorange,mark=square*,thick,error bars/.cd,y dir=both,y explicit] table[x=p,y=value,y error minus=low,y error plus=high,col sep=comma]{figures/sweep.generated.csv};
\end{axis}\end{tikzpicture}
\caption{The simulated population loses full credit as reads disappear. Each point averages \NSeeds{} simulated seed-level suites. Bars show \CiLevel\% bootstrap intervals over seeds; they do not estimate trained-agent performance or new-task uncertainty.}\label{fig:sweep}\end{figure}

At the middle omission setting, simulated full-credit frequency is \MidSweep{}. Table~\ref{tab:sweep} reports the complete sweep. The endpoints reproduce the reference and no-read regimes. Interior points show a graded loss of full credit rather than a change confined to a single threshold. The mean score stays above the full-credit fraction because some damaged traces retain partial credit.\par
\begin{table}[ht]\centering\small\begin{tabular}{lrr}\toprule
Simulated omission probability & Mean score [CI] & Full credit, \% [CI]\\ \midrule
0 & 1.000 [1.000, 1.000] & 100.0 [100.0, 100.0]\\
0.25 & 0.806 [0.773, 0.835] & 78.9 [75.4, 82.0]\\
0.5 & 0.574 [0.544, 0.603] & 53.8 [50.6, 56.9]\\
0.75 & 0.379 [0.355, 0.401] & 33.9 [31.5, 36.1]\\
1 & 0.216 [0.216, 0.216] & 17.5 [17.5, 17.5]\\
\bottomrule\end{tabular}
\caption{Additional simulated population experiment. Intervals bootstrap \NSeeds{} complete seed-level suites at each omission setting. The endpoints are deterministic under this simulated mechanism.}\label{tab:sweep}\end{table}
The simulated curve should not be interpreted as a forecast of agent reliability. It imposes independent read omissions on an otherwise correct, privileged trace. Real policies may compensate for a missed read, omit related actions together, or alter their final report. Such adaptations can change the curve. The experiment instead shows how one precisely defined control loss propagates through the official score.\par
\section{Roadmap}
A useful next study would keep the fault schedule fixed while varying learned agents under the same visible instructions. It should report tool errors separately from provider failures and retain the full trace. A free-only route can support such work, but no hosted inference route was configured or tested for this study. Availability elsewhere remains unknown.\par
A stronger generalization design would introduce held-out task instances and unseen compositions of faults. A transient server error followed by schema drift would test whether a controller can revise its recovery strategy. Status reads that lag writes would test whether verification itself can become ambiguous. These extensions require new expected-state contracts before any scores can be compared with the present version.\par
Fault-free matched tasks would answer a different question. Pairing each fault-bearing task with an otherwise identical clean task could isolate the burden of recovery. Domain experts could then assess whether the fault descriptions and acceptable stopping rules resemble operational practice. Neither comparison appears in the present data. We propose them as validation work, not as completed evidence.\par
The scoring contract also needs explicit review when transferred to other services. Some APIs welcome redundant idempotent writes, while this benchmark rejects writes after known completion. Some services cannot expose a reliable status query. A future version should preserve these distinctions in named task types instead of silently changing the interpretation of an existing score.\par
\section{Caveats}
\paragraph{Internal validity.} The benchmark author supplies the environment and the reference policy. A positive reference score is necessary but not independent validation. The response-aware controller offers a separate decision path, yet it targets the same known fault vocabulary. Its success does not establish resistance to gaming. Replaying traces checks consistency with the simulator, not the correctness of the simulator's assumptions.\par
\paragraph{External validity.} The tasks are synthetic, public, and small in number. A shared recovery contract underlies all domains. There is no human baseline, field study, expert annotation campaign, or deployment test. The simulated population has no fitted relationship to observed agents. Our evidence supports measurement behavior under these conditions and cannot support a claim of production safety.\par
\paragraph{Statistical scope.} Wilson intervals treat task outcomes as if they came from a sampling process. Repeated templates can weaken that interpretation. Fault-stratified bootstrap intervals condition on the fixed mixture and can collapse at boundaries. Simulated seed intervals answer a different question and must not be compared as if they measured the same uncertainty. No multiplicity-adjusted significance claims are made.\par
\paragraph{Reporting and semantics.} The grader checks a structured outcome status and a limited failure-wording pattern. It does not assess every nuance of a user-facing explanation. Its payload normalization accepts some alternate forms while rejecting others. Status reads are trusted and the clock is virtual. These choices improve control but omit failures in the observation channel itself.\par
\paragraph{Contamination and access.} Public fixtures and reference policies make memorization possible. This study does not measure contamination or claim a hidden evaluation. Hosted models were not evaluated, so no provider comparison or free-model availability claim follows. The local study plan was not independently registered. The work remains an unpublished preprint with no venue acceptance claim.\par
\paragraph{Responsible use.} The environments simulate effects and do not issue real invoices, messages, calendar events, or stock changes. Scores should guide investigation of recovery behavior, not certify authority to act on a user's behalf. A low score identifies a contract failure in this setting. A high score leaves service-specific hazards and human oversight requirements unresolved.\par
\section{Citation and availability}
The source, study plan, generated data, and reproduction instructions accompany the benchmark at \url{https://github.com/pooriaarab/benchmark-tool-errors}. The package remains MIT licensed. This paper and its generated figures use CC-BY-4.0. Cite Pooria Arab, \emph{\PaperTitle}, unpublished preprint. The software citation remains separate from the paper citation.\par
\section{Authors}
Pooria Arab is an independent researcher and the benchmark author. AI assistance supported experiment implementation, analysis, and manuscript drafting. The author remains responsible for the claims and release decision. No independent human review or participant study is claimed. The present packet is prepared for review and does not represent publication in a venue.\par
\bibliographystyle{plainnat}
\bibliography{references}
\end{document}
